%!TEX root = ../thesis.tex
%*******************************************************************************
%****************************** Third Chapter **********************************
%*******************************************************************************
\chapter{Bridging Theory and Scientific Practice in a Data-Driven World}

% **************************** Define Graphics Path **************************
\ifpdf
    \graphicspath{{content/chapter3/figs/raster/}{content/chapter3/figs/pdf/}{content/chapter3/figs/}}
\else
    \graphicspath{{content/chapter3/figs/vector/}{content/chapter3/figs/}}
\fi

From statistical models to applications: case studies and technology, Impact of cloud computing, AI, and machine learning.
\section{Key Challenges}
Reproducibility: Ensuring results can be replicated.
Access to Methods: Democratizing advanced tools.
Transparency: Building trust through open methods.
FAIR Principles


Scientific software differs from traditional software engineering by prioritizing scientific goals, involving developers as domain experts, embracing agile-like and open source practices, and focusing on continuous evolution and unique verification methods, such as bit-level reproducibility, to support complex, evolving scientific inquiries (Easterbrook and Johns).

\section{Current Solutions and Trends}

Why Xarray and Dask?
The Pangeo project strongly encourages the use of Xarray data structures wherever possible. Xarray Dataset and DataArrays contain multidimensional numeric array data and also the metadata describing the data’s coordinates, labels, units, and other relevant attributes. Xarray makes it easy to keep this important metadata together with the raw data; applications can then take advantage of the metadata to perform calculations or create visualizations in a coordinate-aware fashion. The use of Xarray eliminates many common bugs, reduces the need to write boilerplate code, makes code easier to understand, and generally makes users and developers happier and more productive in their day-to-day scientific computing.

Xarray’s data model is explicitly based on the CF Conventions, a well-established community standard which encompasses many different common scenarios encountered in Earth System science. However, Xarray is flexible and does not require compliance with CF conventions. We encourage Pangeo packages to follow CF conventions wherever it makes sense to do so.

Most geoscientists have encountered the CF data model via the ubiquitous netCDF file format. While Xarray can easily read and write netCDF files, it doesn’t have to. This is a key difference between software built on Xarray and numerous other tools designed to process netCDF data (e.g. nco, cdo, etc. etc.): Xarray data can be passed directly between python libraries (or over a network) without ever touching a disk drive. This “in-memory” capability is a key ingredient to the big data scalability of Pangeo packages. Very frequently the bottleneck in data processing pipelines is reading and writing files.

Another important aspect of scalability is the use of Dask for parallel and out-of-core computations. The raw data underlying Xarray objects can be either standard in-memory numpy arrays or Dask arrays. Dask arrays behave nearly identically to numpy arrays (they support the same API), but instead of storing raw data, they store a symbolic computational graph of operations. An example computational graph would start by reading data from disk or network and then preform transformations or mathematical calculations. No operations are actually executed until actual numerical values are required, such as for making a figure. This is called lazy execution. Dask figures out how to execute these computational graphs efficiently on different computer architectures using sophisticated techniques. By chaining operations on dask arrays together, researchers can symbolically represent large and complex data analysis pipelines and then deploy them effectively on large computer clusters.


Welcome to pyOpenSci
We support the scientific Python tools that drive open science through peer review, training and community building.

\paragraph{Data Repositories and Archives}

Earth System Grid Federation (ESGF)
Overview: ESGF is a global infrastructure that provides access to a vast array of climate data, including simulations from CMIP, observational datasets, and reanalysis products.
Key Features: Distributed data storage, federated data access, and tools for data search, retrieval, and analysis.


Copernicus Climate Change Service (C3S)
Overview: C3S is part of the European Union's Copernicus program, providing information about the past, present, and future climate. It offers access to a wealth of climate data and analysis tools.
Key Features: Climate data store, reanalysis datasets, and climate indicators.

Climate Data Store (CDS)
Overview: Managed by the European Centre for Medium-Range Weather Forecasts (ECMWF), CDS offers a wealth of climate data, including satellite observations, in situ data, and climate reanalysis.
Key Features: Comprehensive data catalog, tools for data access and processing, APIs for integration.

Copernicus Atmosphere Monitoring Service (CAMS)
Overview: CAMS provides comprehensive data on atmospheric composition, including greenhouse gases, aerosols, and reactive gases. It supports climate monitoring and air quality forecasting.
Key Features: Global data, forecasting capabilities, and detailed emissions inventories.

\paragraph{Modelling and Simulation Platforms}

CMIP (Coupled Model Intercomparison Project)
Overview: CMIP coordinates climate model experiments that compare and analyze the outputs of various climate models from around the world. It plays a crucial role in assessing climate projections and understanding model uncertainties.
Key Features: Standardized data formats, community-developed protocols, and a central repository for model outputs.

Google Earth Engine
Overview: A cloud-based platform for planetary-scale environmental data analysis. It provides access to a vast public data archive and powerful computing resources to analyze geospatial data.
Key Features: Scalable cloud computing, extensive data catalog, APIs for custom analysis.

OpenClimateGIS
Overview: An open-source toolkit for accessing, analyzing, and visualizing climate data. It integrates with GIS tools to provide spatial analysis capabilities.
Key Features: Python-based, support for various data formats, integration with GIS systems.

JupyterHub for Science
Overview: JupyterHub is an open-source platform that allows researchers to share computational resources, code, and data within a collaborative environment. It supports interactive computing and reproducibility.
Key Features: Multi-user Jupyter notebooks, support for various programming languages, integration with cloud and HPC resources.

Open Science Movement, FAIR, open access, open data, open-source software.

Impact: Transforming research collaboration and dissemination.
Python Libraries:
\begin{itemize}
    \item Xarray: Multi-dimensional arrays.
    \item Dask: Scalable computing.
    \item Zarr: Efficient storage.
\end{itemize}
Innovative Tools: Docker, Snakemake, Apache Airflow, Git/GitHub.

other relevant open source software

\section{Introduction to the xeofs Library}
Overview: Design and key features
Core Functionalities: PCA techniques and variants;
Scalable computations; Integration with Python libraries;
Pre- and post-processing; Custom methods and bootstrapping

ensuring reproducibility and transparency through extensive collaborative coding, documentation and version control.
Examples of reproducible research with xeofs. (all of our studies + other studies that have used xeofs)


open source software helps 
\begin{itemize}
    \item to democratize access to cutting-edge tools and techniques, enabling researchers to focus on their scientific questions rather than technical challenges.
    \item to improve reproducibility and transparency in research by providing open-source code that can be easily shared and reused by the community.
\end{itemize}


xeofs addresses this through the following tools and practices for reproducible Science
\begin{itemize}
    \item Version Control Systems: Discuss the use of version control systems like Git/Github for tracking changes in code and data.
    \item Open Data Practices: Explain the importance of open data repositories and how they facilitate reproducibility (e.g., Zenodo, Figshare).
    \item Documentation Standards: Emphasize the need for thorough documentation of methods, data, and code (readthedocs, sphinx, Jupyter Notebooks).
\end{itemize}
also address F.A.I.R principles


in particular xeofs allows for the following features:
\begin{enumerate}
    \item Streamlined Data Analysis for various PCA techniques and variants: xeofs facilitates the application of dimensionality reduction techniques such as EOF analysis (commonly referred to as Principal Component Analysis (PCA) in other fields). This is crucial for extracting meaningful patterns from large, multi-dimensional datasets commonly encountered in climate science.
    \item Scalable computations with dask and zarr: efficient scaling of computations across multiple cores or clusters; allows for out-of-memory processing and improved performance thanks to randomized SVD; essential for handling the extensive datasets typical in climate research
    \item Seamless Integration with state-of-the-art Python libaries: xarray, dask, and zarr; ensures compatibility with existing workflows and tools, simplifying the analysis process, scikit-learn interface
    \item streamlines pre and postprocessing steps like normalization, weighting and handling of missing values
    \item bootstrapping module for PCA model evaluation
    \item Custom dimensionality reduction methods: allows for the introduction of custom dimensionality reduction methods, crucial for researchers who require specialized analysis techniques tailored to their unique data and research questions
\end{enumerate}

\mynote{performance analysis of full SVD compared to randomized SVD; xeofs vs. eofs; reference back to Theory PCA}


